% #############################################################################
% This is Chapter 1
% !TEX root = ../main.tex
% #############################################################################
% Change the Name of the Chapter i the following line
\fancychapter{Introduction}
% The following line allows to ref this chapter
\label{chap:intro}

Type systems are a powerful tool which can be used to ensure the absence of some runtime errors, which would otherwise require extensive testing or program analysis. Essentially, type systems form a set of rules which limit the programs that can be written in a given language. Most traditional compiled imperative languages, such as C and Java, have type systems that avoid some common programming errors, such as passing a string to a function that expects an integer. Although such programs in a weakly or dynamically typed language such as Python or JavaScript would compile and run, they would likely fail at runtime, whereas in a strongly typed language, the compiler would catch the error at compile time. Strongly typed languages can additionally make use of their type to provide more information about the behavior of the code to the compiler, opening up new optimization paths.

Some languages take this concept further, by using the type system to enforce stronger properties of the code. For example, Rust \cite{Matsakis:2014} uses an affine type system to enforce memory safety, absence of data races and dangling pointers, which are common source of bugs in C and C++ programs, without the need for a garbage collector. However, Rust still does not guarantee that concurrent programs are deadlock-free, as it does not have a type system that enforces communication protocols.

% \todo{Maybe also mention here that type-safe languages are becoming more and more popular and valuable, e.g., US government encouraging writing new software in Rust instead of C/C++ due to memory vulnerabilities}.

Session types \cite{Honda:1993,Honda:1998} have been used to enforce correctness in communication protocols, namely in languages based on process calculus, which model concurrent systems \cite{Milner:1980,Milner:1992:1,Milner:1992:2}. By coupling session types with linear logic \cite{Girard:1987,Caires:2010}, it has been possible to additionally guarantee deadlock freedom. \ac{CLASS} \cite{Rocha:2022}, is an example of a session-typed language, built on top of the correspondence between session types and linear logic, that provides these guarantees. In fact, \ac{CLASS} is the first language to simultaneously guarantee memory safety and the absence of dead-locks. The original \ac{CLASS} interpreter implemented by Rocha et al. \cite{Rocha:2021:artifact} served as a proof of concept for the language. It wasn't optimized for performance, and thus, it wasn't suitable for real-world applications.

The Linear Session Abstract Machine (\ac{SAM}) \cite{Caires:2024} is an abstract machine that allows for the sequential execution of session-typed languages such as \ac{CLASS} through the use of buffered communication channels. Caires et al. expanded the existing \ac{CLASS} interpreter with this new scheme \cite{Caires:2024:artifact}. This permitted a great increase in performance, as the overhead of spawning threads is very high in modern systems, and thus it is not feasible to spawn threads for every parallel process in a session-typed program, as most threads end up executing few instructions and spend most time waiting for communication. However, the \ac{SAM} still does not make use of the properties of the source language to generate efficient machine code, as it is a high-level abstract machine.

During this work, a compilation scheme was designed and implemented which translates \ac{CLASS} programs into efficient C code, while maintaining all of the safety and liveness guarantees of \ac{CLASS}. This approach allows the execution of \ac{CLASS} programs in a orders of magnitude more efficient manner than the \ac{SAM} interpreter, as the compiler can make use of the extensive knowledge of the compiled programs. The \ac{SAM} interpreter, on which the compiler is based, currently has no support for the affine sessions, reference cells and shared state features of \ac{CLASS}. Instead, it covers only the pure fragment of the language. The developed compiler supports all of the features of \ac{CLASS} in the compiled programs, including the aforementioned ones.

The base \ac{CLASS} language was also slightly extended with a program abort construct, and a input extraction construct, allowing programs to read input from the standard input. This allowed for the creation of more complex programs, which were helpful to measure the real-world performance of the compiled programs. The performance of the compiled programs was compared with the original \ac{CLASS} interpreter, the \ac{SAM} interpreter, and with equivalent programs written in Haskell. The results show that the compiled programs are orders of magnitude faster than the interpreted ones, and, in some cases, outperform equivalent Haskell programs compiled with GHC.

\todo{Review/rewrite}

\todo{Brief description of the structure of the document}
